\documentclass[tikz,border=8pt]{standalone}
\usepackage{fontspec}
\usepackage{xeCJK}
\IfFontExistsTF{Noto Sans TC}{\setCJKmainfont{Noto Sans TC}}{\setCJKmainfont{Microsoft JhengHei}}
\setmainfont{Arial}
\usetikzlibrary{arrows.meta,positioning,calc}
\definecolor{paper}{HTML}{FBFAF7}
\definecolor{navy}{HTML}{152238}
\definecolor{muted}{HTML}{647184}
\definecolor{blue}{HTML}{9AB6FF}
\definecolor{bluesoft}{HTML}{EEF3FF}
\definecolor{green}{HTML}{2D7A55}
\definecolor{greensoft}{HTML}{EDF7F2}
\definecolor{gold}{HTML}{B38B2E}
\definecolor{goldsoft}{HTML}{FFF8E6}
\definecolor{line}{HTML}{D9DDE5}
\begin{document}
\begin{tikzpicture}[x=1cm,y=1cm,>=Latex]
\path[fill=paper] (0,0) rectangle (16,9.6);
\node[font=\bfseries\fontsize{18}{22}\selectfont,text=navy] at (8,8.82) {Vibe Coding 要先寫成需求合約};
\node[font=\fontsize{10.5}{13}\selectfont,text=muted] at (8,8.12) {中文需求不是咒語，要拆成資料、例外與查核};

\draw[rounded corners=9pt,fill=goldsoft,draw=gold,line width=.8pt] (.95,4.0) rectangle (4.4,5.55);
\node[font=\bfseries\fontsize{11}{14}\selectfont,text=navy] at (2.68,5.03) {白話需求};
\node[font=\fontsize{8.5}{11}\selectfont,text=muted] at (2.68,4.62) {幫我抓財報資料};

\draw[rounded corners=9pt,fill=bluesoft,draw=blue,line width=.8pt] (6.0,6.0) rectangle (9.7,7.15);
\node[font=\bfseries\fontsize{10.5}{13}\selectfont,text=navy] at (7.85,6.72) {資料定義};
\node[font=\fontsize{8}{10}\selectfont,text=muted] at (7.85,6.38) {來源、欄位、期間、單位};
\draw[rounded corners=9pt,fill=white,draw=line,line width=.8pt] (6.0,4.25) rectangle (9.7,5.4);
\node[font=\bfseries\fontsize{10.5}{13}\selectfont,text=navy] at (7.85,4.97) {例外處理};
\node[font=\fontsize{8}{10}\selectfont,text=muted] at (7.85,4.63) {缺值、格式錯、重複檔};
\draw[rounded corners=9pt,fill=greensoft,draw=green,line width=.8pt] (6.0,2.5) rectangle (9.7,3.65);
\node[font=\bfseries\fontsize{10.5}{13}\selectfont,text=navy] at (7.85,3.22) {查核規則};
\node[font=\fontsize{8}{10}\selectfont,text=muted] at (7.85,2.88) {合計、期間、人工抽查};

\draw[rounded corners=9pt,fill=navy,draw=navy,line width=.8pt] (11.55,3.68) rectangle (15.05,5.85);
\node[font=\bfseries\fontsize{12}{15}\selectfont,text=white] at (13.3,5.05) {會計洞察};
\node[font=\fontsize{8.5}{11}\selectfont,text=white] at (13.3,4.62) {程式最後要回答問題};

\draw[->,line width=1pt,draw=navy] (4.48,4.78)--(5.82,6.58);
\draw[->,line width=1pt,draw=navy] (4.48,4.78)--(5.82,4.82);
\draw[->,line width=1pt,draw=navy] (4.48,4.78)--(5.82,3.08);
\draw[->,line width=1pt,draw=navy] (9.85,6.58)--(11.34,5.15);
\draw[->,line width=1pt,draw=navy] (9.85,4.82)--(11.34,4.82);
\draw[->,line width=1pt,draw=navy] (9.85,3.08)--(11.34,4.48);

\draw[draw=line,line width=.7pt] (.95,1.35)--(15.05,1.35);
\node[font=\fontsize{10.5}{13}\selectfont,text=navy] at (8,.82) {程式能力的入口，是把需求說到足以被檢查};
\end{tikzpicture}
\end{document}
